\subsection{TinyMLOps: MLOps for Resource-Constrained Devices}

The rapid proliferation of Artificial Intelligence on embedded systems has necessitated the adaptation of traditional Machine Learning Operations (MLOps) principles to the specific constraints of micro-controllers. This emerging field, known as \textbf{TinyMLOps}, addresses the lifecycle management of AI models in environments where memory, compute, and power are severely limited.

\subsubsection{Distinction from Cloud-Scale MLOps}

While traditional MLOps focuses on scaling model deployment across cloud clusters, managing petabyte-scale datasets, and orchestrating microservices, TinyMLOps must prioritize high-fidelity hardware awareness. The primary differences are summarized in Table~\ref{tab:mlops_vs_tinymlops}.

\begin{table}[h]
\centering
\caption{Comparison between traditional Cloud MLOps and TinyMLOps.}
\label{tab:mlops_vs_tinymlops}
\begin{tabular}{|l|l|l|}
\hline
\textbf{Feature} & \textbf{Cloud MLOps} & \textbf{TinyMLOps} \\
\hline
Target Hardware & GPU Clusters / Data Centers & MCUs (KB-scale RAM) \\
Optimization & Throughput / Latency & Memory / Energy Efficiency \\
Deployment Unit & Docker Containers / APIs & Firmware Binary Images \\
Connectivity & Persistent / High-bandwidth & Intermittent or Offline \\
Metric Priority & Prediction Accuracy & Accuracy vs. Resource Footprint \\
\hline
\end{tabular}
\end{table}

\subsubsection{Core Pillars of TinyMLOps}

My proposed framework implements three fundamental pillars of the TinyMLOps paradigm:

\begin{itemize}
    \item \textbf{Hardware-Aware Optimization}: Unlike cloud models that can easily scale with more VRAM, embedded models must undergo structural transformations (e.g., quantization, pruning, and "neural surgery") to fit into specific MCU memory profiles.
    
    \item \textbf{Integrated Automation ("Plumbing")}: TinyMLOps requires the seamless connection of traditionally disjoint tools—Python-based training environments, C-based firmware toolchains (like STM32Cube), and binary flashing utilities.
    
    \item \textbf{Reproducibility}: By defining the end-to-end workflow as a declarative graph (e.g., via LangGraph), I ensure that an optimized model can be reconstructed and re-verified consistently across different hardware revisions.
\end{itemize}

By adopting a TinyMLOps approach, the development of embedded AI transitions from a manual, expert-intensive process to a systematic, automated engineering discipline.
